% topic: algorithms
\question คุณครูมีลูกอมอยู่กล่องหนึ่ง ต้องการแจกให้นักเรียนจำนวน \textbf{10 คน} ที่ยืนเรียงแถวกัน โดยมีกฎการแจกดังนี้:

\begin{itemize}
  \item นักเรียนแต่ละคนต้องได้ลูกอม\textbf{อย่างน้อย 1 ชิ้น}
  \item นักเรียนที่มีคะแนน\textbf{มากกว่า}เพื่อนที่อยู่ติดกัน (ซ้ายหรือขวา) จะต้องได้ลูกอม\textbf{มากกว่า}เพื่อนคนนั้น
  \item คุณครูต้องการใช้ลูกอมจำนวน\textbf{น้อยที่สุด}ที่เป็นไปได้
\end{itemize}

คะแนนของนักเรียนทั้ง 10 คนจากซ้ายไปขวา มีดังนี้:

\begin{center}
\begin{tabular}{|c|c|c|c|c|c|c|c|c|c|c|}
\hline
\textbf{คนที่} & 1 & 2 & 3 & 4 & 5 & 6 & 7 & 8 & 9 & 10 \\
\hline
\textbf{คะแนน} & 3 & 5 & 5 & 4 & 2 & 3 & 7 & 7 & 6 & 4 \\
\hline
\end{tabular}
\end{center}

\medskip
\textbf{คำถาม:} คุณครูจะต้องใช้ลูกอม\textbf{อย่างน้อยที่สุด}ทั้งหมดกี่ชิ้น (ตอบเป็นตัวเลขจำนวนเต็ม)

\par\medskip
ตอบ ในกระดาษคำตอบ
% answer: 20
% solution:
%   Standard "candy" problem — two-pass greedy.
%   Scores: [3, 5, 5, 4, 2, 3, 7, 7, 6, 4]
%
%   Pass 1 (left to right):
%   If score[i] > score[i-1]: candy[i] = candy[i-1] + 1
%   Else: candy[i] = 1
%   -> [1, 2, 1, 1, 1, 2, 3, 1, 1, 1]
%
%   Pass 2 (right to left):
%   If score[i] > score[i+1]: candy[i] = max(candy[i], candy[i+1] + 1)
%   -> [1, 2, 3, 2, 1, 2, 3, 3, 2, 1]
%
%   Total = 1+2+3+2+1+2+3+3+2+1 = 20
